\section{Conclusion}
\label{sec:conclusion}

We presented \ARES{}, a compact formal reliability framework for
detecting anomalous behaviour in agentic AI systems deployed in
cybersecurity environments. The conference contribution is narrow:
the SKC mismatch model produces a learnable ARS that generalises
across threat classes (test F1 0.96), and the Epistemic Gap
metric formalises the hallucination precondition as
\textit{confident ignorance} and decomposes false-negative risk into
knowledge and calibration components. A cascade reliability bound for
tightly-coupled pipelines and a temporal-drift model for knowledge
ageing are reported as supporting results.

Practically, ARS supports pre-deployment reliability screening, the
Epistemic Gap supports runtime hallucination-risk monitoring, the
temporal-drift model supports evidence-based knowledge maintenance, and
the cascade bound informs conservative multi-agent design, an initial
formal foundation for trustworthy agentic AI in cybersecurity that may
generalise to other domains where autonomous agents decide under
incomplete knowledge.

Future work follows the limitations above. The most consequential step
is to move from controlled injection to measurement: recovering the
mismatch dimensions from live agents in a production SOC, where
$\delta_K$ and expressed uncertainty must be read from real telemetry
rather than set by construction. Companion work takes a first step for
the knowledge and calibration terms, and early results there suggest
that agent calibration, not knowledge coverage alone, governs whether a
mismatch becomes a dangerous failure. Cyclic multi-agent architectures
and adversarial mismatch injection remain open. The \SAFK{} corpus and
evaluation testbed are released
publicly.\footnote{Project page: \url{https://github.com/IntelEyes/ARES-Framework}}
